\subsection{Estimación del coste de desarrollo}

Para estimar el coste de desarrollo del sistema se ha realizado una planificación basada en las diferentes fases habituales del ciclo de vida del software: análisis de requisitos, diseño del sistema, desarrollo, pruebas y documentación. Este enfoque se corresponde con los modelos clásicos de ingeniería del software descritos en la literatura académica \cite[Chap.~2, pp.~27--37]{sommerville}.

Dado que el proyecto ha sido desarrollado en el contexto de un Trabajo Fin de Grado por un único desarrollador, se ha considerado una estimación de costes basada en el salario medio de un desarrollador web en España. Según datos recopilados por la plataforma de empleo Indeed, el salario medio por hora para este perfil profesional se sitúa aproximadamente en \textbf{15,88~\euro/hora} \cite{indeed2024salarios}.

A partir de esta tarifa media y de la estimación de horas necesarias para cada fase del proyecto, se obtiene el coste total de desarrollo que se muestra en la Tabla~\ref{tabla_coste_desarrollo}.

\begin{table}[H]
\centering
\small
\caption{Estimación de costes de desarrollo del sistema}
\label{tabla_coste_desarrollo}
\begin{tabularx}{\textwidth}{|L{2.8cm}|Y|c|c|c|}
\hline
\textbf{Fase del proyecto} & \textbf{Descripción} & \textbf{Horas} & \textbf{\euro/hora} & \textbf{Coste (\euro)} \\ \hline

Análisis de requisitos &
Reuniones con el cliente, estudio del problema, definición de funcionalidades y documentación inicial &
30 & 15,88 & 476,40 \\ \hline

Diseño de arquitectura &
Definición de arquitectura del sistema, elección de tecnologías y diseño general del sistema &
25 & 15,88 & 397,00 \\ \hline

Diseño de base de datos &
Modelado de datos, diseño de tablas, relaciones y optimización inicial &
20 & 15,88 & 317,60 \\ \hline

Diseño de interfaz (UI/UX) &
Diseño de pantallas, estructura de navegación y experiencia de usuario &
30 & 15,88 & 476,40 \\ \hline

Desarrollo backend &
Implementación de lógica de negocio, APIs, autenticación, gestión de pedidos y usuarios &
90 & 15,88 & 1429,20 \\ \hline

Desarrollo frontend &
Implementación de la interfaz web, integración con APIs y diseño responsive &
100 & 15,88 & 1588,00 \\ \hline

Integración de servicios externos &
Integración con APIs externas como servicios de mapas u otras herramientas &
20 & 15,88 & 317,60 \\ \hline

Pruebas y validación &
Testing funcional, detección de errores y validación del sistema &
40 & 15,88 & 635,20 \\ \hline

Documentación del proyecto &
Elaboración de documentación técnica y memoria del proyecto &
35 & 15,88 & 555,80 \\ \hline

\textbf{Total desarrollo} & & \textbf{390} & & \textbf{6193,20} \\ \hline

\end{tabularx}
\end{table}

\subsection{Costes de infraestructura}

Además del coste de desarrollo, un sistema software requiere una infraestructura mínima para su despliegue y funcionamiento. En el caso de aplicaciones web modernas, esta infraestructura suele estar compuesta por un servidor de alojamiento, una base de datos gestionada, un dominio web y posibles servicios externos o APIs.

Actualmente existen múltiples plataformas cloud que permiten desplegar aplicaciones web con costes reducidos, especialmente para proyectos de pequeña y mediana escala. Servicios como Vercel, Render o Railway ofrecen planes que oscilan entre los 10 y 25 euros mensuales dependiendo del consumo de recursos \cite{vercelpricing2024}.

Por otra parte, las bases de datos gestionadas basadas en PostgreSQL, como Supabase o Neon, permiten operar con costes similares en proyectos de baja carga. Asimismo, el registro de dominios web suele tener un coste aproximado de entre 10 y 15 euros anuales dependiendo del proveedor \cite{namecheap2024domain}.

Teniendo en cuenta estos valores de mercado, se ha realizado una estimación del coste anual de infraestructura necesaria para el funcionamiento del sistema, la cual se presenta en la Tabla~\ref{tabla_coste_infraestructura}.

\begin{table}[H]
\centering
\small
\caption{Costes de infraestructura del sistema}
\label{tabla_coste_infraestructura}
\begin{tabularx}{\textwidth}{|L{4cm}|Y|c|}
\hline
\textbf{Concepto} & \textbf{Descripción} & \textbf{Coste anual (\euro)} \\ \hline

Servidor / Hosting &
Alojamiento del backend y frontend de la aplicación &
180 \\ \hline

Base de datos &
Servicio gestionado de base de datos PostgreSQL &
120 \\ \hline

Dominio web &
Registro y renovación anual del dominio de la aplicación &
12 \\ \hline

APIs externas &
Servicios externos utilizados por la aplicación como APIs de mapas u otros servicios &
150 \\ \hline

\textbf{Total infraestructura anual} & & \textbf{462} \\ \hline

\end{tabularx}
\end{table}

\subsection{Costes de mantenimiento}

Una vez desarrollado el sistema, es necesario considerar los costes asociados al mantenimiento del software. El mantenimiento incluye tareas como la corrección de errores, actualizaciones de seguridad, adaptación a cambios en los servicios externos utilizados por la aplicación y pequeñas mejoras funcionales.

En ingeniería del software es habitual estimar el coste de mantenimiento anual como un porcentaje del coste inicial de desarrollo. Diversos estudios indican que el mantenimiento de software puede representar entre el 10\% y el 20\% del coste de desarrollo anual dependiendo de la complejidad del sistema \cite{pressman2014software}.

Para este proyecto se ha considerado un valor intermedio del \textbf{15\% del coste de desarrollo}, lo cual proporciona una estimación razonable del esfuerzo necesario para mantener el sistema operativo y actualizado. El desglose del coste anual de mantenimiento se presenta en la Tabla~\ref{tabla_coste_mantenimiento}.

\begin{table}[H]
\centering
\small
\caption{Costes de mantenimiento del sistema}
\label{tabla_coste_mantenimiento}
\begin{tabularx}{\textwidth}{|L{4.4cm}|Y|c|}
\hline
\textbf{Concepto} & \textbf{Descripción} & \textbf{Coste anual (\euro)} \\ \hline

Mantenimiento del sistema &
Corrección de errores, actualizaciones, mejoras y soporte técnico &
928,98 \\ \hline

Infraestructura &
Coste anual de servidor, base de datos, dominio y APIs &
462 \\ \hline

\textbf{Coste anual total de operación} & & \textbf{1390,98} \\ \hline

\end{tabularx}
\end{table}

\subsection{Coste total del proyecto}

En base a las estimaciones realizadas, el coste total del proyecto durante su primer año de funcionamiento asciende a aproximadamente \textbf{7.584,18~\euro}. Este coste incluye el desarrollo inicial del sistema, la infraestructura necesaria para su despliegue y el mantenimiento del software. A partir del segundo año, el coste de operación se reduce a aproximadamente 1.390,98~\euro anuales, correspondientes principalmente al mantenimiento del sistema y a los costes de infraestructura. El desglose completo de los costes se resume en la Tabla~\ref{tabla_coste_total}.

\begin{table}[H]
\centering
\small
\caption{Resumen económico del proyecto}
\label{tabla_coste_total}
\begin{tabularx}{\textwidth}{|Y|c|}
\hline
\textbf{Concepto} & \textbf{Coste (\euro)} \\ \hline

Desarrollo inicial del sistema & 6193,20 \\ \hline

Infraestructura primer año & 462 \\ \hline

Mantenimiento primer año & 928,98 \\ \hline

\textbf{Coste total estimado primer año} & \textbf{7584,18} \\ \hline

\end{tabularx}
\end{table}
